\section{互信息：不确定性的减少量}
\warmth{0}\quad\whisper{互信息把“收到一个信源后，对另一个信源多了解了多少”变成公式。}

\begin{strand}
如果 \(X\) 是信源，\(Y\) 是信宿或接收端观察到的量，那么互信息 \(\info(X;Y)\)
衡量的是：观察 \(Y\) 以后，关于 \(X\) 的不确定性平均减少了多少。
\end{strand}

\block{定义与等价写法}
\begin{definition}[互信息]
互信息定义为
\[
\info(X;Y)=H(X)-H(X|Y).
\]
等价地，也可以写成
\[
\info(X;Y)=H(Y)-H(Y|X)
=H(X)+H(Y)-H(X,Y).
\]
因此，互信息可以看成 \(X\) 与 \(Y\) 共享的那部分信息。
\end{definition}

\begin{proposition}[概率形式]
互信息也可以写成联合分布与边缘分布的比值：
\[
\info(X;Y)
=\sum_{i=1}^{n}\sum_{j=1}^{m}
\pmf(x_i,y_j)
\log\frac{\pmf(x_i,y_j)}{\pmf(x_i)\pmf(y_j)}.
\]
等价地，
\[
\info(X;Y)
=\mathbb{E}\left[\log\frac{\pmf(X|Y)}{\pmf(X)}\right]
=\mathbb{E}\left[\log\frac{\pmf(Y|X)}{\pmf(Y)}\right].
\]
\end{proposition}

\begin{proof}
从
\[
\info(X;Y)=H(X)-H(X|Y)
\]
出发。因为
\[
\pmf(x_i)=\sum_{j=1}^{m}\pmf(x_i,y_j),
\]
所以
\[
\begin{aligned}
H(X)
&=-\sum_{i=1}^{n}\pmf(x_i)\log \pmf(x_i)\\
&=-\sum_{i=1}^{n}\sum_{j=1}^{m}
\pmf(x_i,y_j)\log \pmf(x_i).
\end{aligned}
\]
另一方面，
\[
H(X|Y)
=-\sum_{i=1}^{n}\sum_{j=1}^{m}
\pmf(x_i,y_j)\log \pmf(x_i|y_j).
\]
两式相减，得到
\[
\begin{aligned}
\info(X;Y)
&=\sum_{i=1}^{n}\sum_{j=1}^{m}
\pmf(x_i,y_j)
\left[\log \pmf(x_i|y_j)-\log \pmf(x_i)\right]\\
&=\sum_{i=1}^{n}\sum_{j=1}^{m}
\pmf(x_i,y_j)
\log\frac{\pmf(x_i|y_j)}{\pmf(x_i)}.
\end{aligned}
\]
再由条件概率
\[
\pmf(x_i|y_j)=\frac{\pmf(x_i,y_j)}{\pmf(y_j)}
\]
可得
\[
\info(X;Y)
=\sum_{i=1}^{n}\sum_{j=1}^{m}
\pmf(x_i,y_j)
\log\frac{\pmf(x_i,y_j)}{\pmf(x_i)\pmf(y_j)}.
\]
这也说明
\[
\info(X;Y)
=\mathbb{E}\left[\log\frac{\pmf(X|Y)}{\pmf(X)}\right].
\]
同理从 \(\info(X;Y)=H(Y)-H(Y|X)\) 出发，可以得到
\[
\info(X;Y)
=\mathbb{E}\left[\log\frac{\pmf(Y|X)}{\pmf(Y)}\right].
\]
\end{proof}

\begin{yourturn}
复述互信息等价式证明的关键动作：
\[
H(X)=-\sum_i \pmf(x_i)\log\pmf(x_i)
=-\sum_{i,j}\answerin[2.8cm]{\pmf(x_i,y_j)}\log\pmf(x_i),
\]
所以 \(H(X)-H(X|Y)\) 可以合并成同一个 \((i,j)\) 双重求和。
\workspace[2]
\end{yourturn}

\block{性质}
\noindent\begin{tabularx}{\linewidth}{@{}l L@{}}
\toprule
{\color{indigo}性质} & \multicolumn{1}{l}{\color{indigo}含义}\\
\midrule
非负性 & \(\info(X;Y)\geq 0\)。平均来看，观察不会让不确定性“负减少”。\\
对称性 & \(\info(X;Y)=\info(Y;X)\)。\(Y\) 中关于 \(X\) 的信息量，等于 \(X\) 中关于 \(Y\) 的信息量。\\
独立时为零 & 若 \(X\) 与 \(Y\) 独立，则 \(\info(X;Y)=0\)。\\
完全确定时最大 & 若 \(Y\) 能完全确定 \(X\)，则 \(\info(X;Y)=H(X)\)。\\
\bottomrule
\end{tabularx}

\begin{remark}
在固定信道转移概率 \(\pmf(y|x)\) 时，\(\info(X;Y)\) 是输入分布 \(\pmf(x)\) 的凹函数；
在固定输入分布 \(\pmf(x)\) 时，\(\info(X;Y)\) 是信道转移概率 \(\pmf(y|x)\) 的凸函数。
这解释了两个后续优化问题：固定信道 \(\pmf(y|x)\) 时，最大化 \(\info(X;Y)\) 得信道容量；
固定失真约束时，最小化 \(\info(X;\hat X)\) 得率失真函数 \(R(D)\)。
\end{remark}

\begin{yourturn}
用熵的图像直觉填空：
\[
\info(X;Y)=H(X)-H(X|Y)
\]
表示观察 \(Y\) 以后，关于 \(X\) 的不确定性减少了
\answerin[3cm]{\info(X;Y)}。
\workspace[2]
\end{yourturn}

\begin{yourturn}
若 \(X\) 与 \(Y\) 独立，则 \(\pmf(x_i,y_j)=\pmf(x_i)\pmf(y_j)\)。代入概率形式可得
\[
\info(X;Y)=\sum_{i,j}\pmf(x_i,y_j)\log\answerin[1.5cm]{1}=\answerin[1.5cm]{0}.
\]
\workspace[2]
\end{yourturn}

\loose{下一步接信源编码：无失真看 \(H(X)\)，限失真看
\(R(D)=\min \info(X;\hat X)\)。}
